\section{Related Work and a Survey of Route Matchers}\label{sec:related}

\subsection{Tries and Lookup Structures}

A route table is a set of keys made of path segments, so the trie is its natural structure.
De~la~Briandais and Fredkin described the digital search tree, which Fredkin named
``trie''~\cite{delabriandais1959file,fredkin1960trie}. PATRICIA compresses a binary trie by skipping
nodes with a single branch~\cite{morrison1968patricia}. The radix trees of several routers follow
this line. The double-array trie encodes a trie in flat arrays~\cite{aoe1989doublearray}. The burst
trie mixes trie nodes with other containers~\cite{heinz2002burst}. The adaptive radix tree chooses
a node's layout by its number of children, and its authors report lookups comparable to hash
tables~\cite{leis2013art}. v2 borrows two of these ideas. It flattens its trie into arrays after the
build, without the double-array encoding. It also changes how a node finds a literal child once the
node has many children (Section~\ref{sec:design}).

Address lookup in routers of the network layer is a large neighbouring field. Its surveys cover
the design space~\cite{ruizsanchez2001survey,gupta2001classification,taylor2005survey}. Its methods
include adaptive branching~\cite{andersson1993adaptive}, level compression~\cite{nilsson1999lctrie},
prefix expansion~\cite{srinivasan1999cpe} and tables small enough for a cache~\cite{degermark1997small}.
Others are binary search over hash tables~\cite{waldvogel1997scalable}, the tree
bitmap~\cite{eatherton2004treebitmap} and Poptrie~\cite{asai2015poptrie}. Poptrie's authors evaluate their structure with an analysis of
processor cycles. We follow that practice and report hardware counters beside every timing.
Name lookup in named data networking is the closest analogue, because names there are hierarchies
of components~\cite{wang2012nce,so2013ndn,yuan2015lnpm}. Those structures find the longest matching
prefix. A route matcher instead needs an exact match with typed parameters, a catch-all, precedence
and captured values.

v2's exact-match table and its hashed trie nodes use open addressing~\cite{peterson1957addressing}.
Robin Hood hashing~\cite{celis1985robinhood}, cuckoo hashing~\cite{pagh2004cuckoo}, and perfect and
minimal perfect hashing~\cite{fredman1984sparse,czech1992optimal,belazzougui2009chd,pibiri2021pthash}
are alternatives.
GNU gperf generates perfect hash functions as source code~\cite{gnu2025gperf}, and the frozen library
offers constant-evaluated hash containers~\cite{guelton2024frozen}. A table that is built when the
program compiles knows its keys, so a perfect hash would be possible there. v2 builds the same layout
at compile time and at run time instead, so that one lookup serves both. We did not test a perfect
hash, and we make no claim about it.

\subsection{Research on HTTP Routing}

Zhang and Xu replace a table of regular expressions with a trie searched depth first, with
static segments in a hash table~\cite{zhang2015urlrouting}. A work-in-progress paper embeds a hash
trie for REST URLs in a server~\cite{dsouza2018htu}. Zeng and colleagues match literal URLs from a
fixed start in HTTP streams, without parameters or precedence~\cite{zeng2017mh}. The Go benchmark
suite of Schmidt compares Go routers on the route sets of real APIs~\cite{schmidt2020gohttproutingbench}.
Our exploratory real-world table is its GitHub API list. That suite times Go routers inside Go.
matchit's documentation ranks Rust routers by a benchmark~\cite{ahmed2026matchit}, and wayfind keeps
a suite that times Rust routers~\cite{sworn2026wayfind}. A gateway benchmark times API gateways on the
route sets of real interfaces~\cite{vm0012024gateways}. Each times routers of one language or one
kind. We time routers of four languages behind one harness, with null arms and checked answers. The path
grammar that a parameter accepts comes from RFC~3986~\cite{bernerslee2005rfc3986}. The meaning of
HEAD and of a 405 answer comes from RFC~9110~\cite{fielding2022rfc9110}.

\subsection{The Routers Measured}\label{sec:survey}

Table~\ref{tab:survey} lists the sixteen competitors. Each was the latest release of its project
when the hypotheses were frozen, or, for r3, the head of its maintained branch. The table says how
each stores its routes, as its source or documentation at that version states it. It also says what
its lookup returns. Appendix~\ref{app:lookup} gives the exact call that the harness times for each.

\begin{table}[htbp]
\caption{The sixteen routers compared with v2. ``Structure'' is read from the source or the
documentation at the version measured. ``Values'' is what the lookup returns for a parameter.
Appendix~\ref{app:config} gives the build of each.\label{tab:survey}}
\begin{adjustwidth}{-\extralength}{0cm}
\centering\small
\begin{tabularx}{\fulllength}{@{}l l l X X@{}}
\toprule
\textbf{Router} & \textbf{Lang.} & \textbf{Version} & \textbf{Structure} & \textbf{Values} \\
\midrule
Crow~\cite{crow2026crow} & C++ & \VerCrow & a trie; each node holds a key, a parameter type and a vector of children & copies into strings \\
Drogon~\cite{drogon2026drogon} & C++ & \VerDrogon & a lower-cased copy of the path in two hash maps, then a regular expression per parameter route, in turn & a new vector of strings \\
Oat++~\cite{oatpp2024oatpp} & C++ & \VerOatpp & a list of patterns, tried in order & views; the tail copied \\
Pistache~\cite{pistache2024pistache} & C++ & \VerPistache & a tree of segments whose children sit in hash maps & copies into strings \\
Boost.URL~\cite{boosturl2026router} & C++ & \VerBoostUrl & the example router of the library: segment templates, literal first & views \\
r3~\cite{r32026r3} & C & \VerRThree & a prefix trie compiled from the paths, with PCRE2 \VerPcre & views \\
uWebSockets~\cite{uwebsockets2026uwebsockets} & C++ & \VerUWebSockets & a tree of segments under the method; children sorted literal, parameter, catch-all & views; none for a catch-all \\
glaze~\cite{berry2026glaze} & C++ & \VerGlaze & a hash map of static paths, then a trie with hashed literal children & decoded strings in a map \\
cpp-httplib~\cite{hirose2026httplib} & C++ & \VerCppHttplib & a list per method, tried in order; a regular expression where a pattern needs one & copies, or views of the expression's groups \\
matchit~\cite{ahmed2026matchit} & Rust & \VerMatchit & a radix trie; a literal before a parameter & offsets \\
actix-router~\cite{actix2026router} & Rust & \VerActixRouter & a list of resource definitions, tried in order & offsets \\
path-tree~\cite{viz2025pathtree} & Rust & \VerPathTree & a compressed prefix tree with literal and parameter children & offsets \\
httprouter~\cite{schmidt2019httprouter} & Go & \VerHttprouter & a compressing radix tree; a literal and a parameter never share a segment & offsets \\
gin~\cite{gin2026gin} & Go & \VerGin & httprouter's tree, extended to let a literal stand beside a parameter & offsets \\
net/http~\cite{go2026go1271,amsterdam2024routing} & Go & Go \VerNetHttp & the ServeMux; the most specific pattern wins & copies \\
chi~\cite{chi2026chi} & Go & \VerChi & a radix tree with literal, regular-expression, parameter and catch-all nodes & offsets \\
\bottomrule
\end{tabularx}
\end{adjustwidth}
\end{table}

Two of these are libraries that larger frameworks build on: matchit, which axum uses, and
actix-router, the router of actix-web. axum's latest release, \VerAxum, pins matchit
\VerAxumMatchit~\cite{axum2026axum}, so the arm measures matchit's latest release rather than the one
that axum uses. gin's tree derives from httprouter's and carries its notice. Drogon's router
needs a running application, so its lookup is transcribed into the harness from its source at the
pinned tag. The transcription is a threat to validity (Section~\ref{sec:threats}).

The routers disagree on what a route means. httprouter refuses a literal and a parameter at the same
position for one method. Go's ServeMux panics at registration when two patterns overlap and neither
is more specific. matchit refuses ambiguous routes when they are inserted. Lists tried in order give
precedence to the earlier route. Section~\ref{sec:results} runs eight probes of the path rules on
every arm and shows where the measured routers differ.

\subsection{Routes Checked When the Program Compiles}

Typed web frameworks check routes against handlers in several languages. In Servant, a web API is a
Haskell type, interpreted as a server, a client and documentation~\cite{mestanogullari2015servant}.
Ocsigen uses OCaml's types for valid links and service calls~\cite{balat2006ocsigen}. Yesod turns a
route language into Haskell code with Template Haskell. Its route quasiquoter stops the compilation
when a route line is malformed or two routes may overlap. Each route's arguments take the types of
its dynamic pieces~\cite{yesod2026routing,yesod2025core}. The mechanism is quasiquotation, which checks a
language written as a string when the program compiles~\cite{mainland2007quasiquoting}. Play
compiles its route file but resolves overlaps by declaration order~\cite{play2026routing}. Rocket
requires each dynamic parameter to appear as a typed argument of its handler. Colliding routes stop
Rocket at launch, not at compile time~\cite{rocket2024rocket}. Beyond routing, Ur/Web is a statically
typed language for web applications~\cite{chlipala2015urweb}.

Two C++ routers among the related work check part of a route when the program compiles. Crow's
route macro derives the argument types from the parameter tags of a pattern, and a handler that
cannot take them does not compile. A pattern without a leading slash, and a route defined twice, are
found at run time~\cite{crow2026routing}. RESTinio builds a route from parsers whose results reach
the handler as separate arguments, so a mismatched handler does not compile. Its route is a C++
expression rather than a pattern string~\cite{stiffstream2025restinio}.

v2's front end builds on the C++ techniques that check a string argument when the program compiles.
The control string of printf decides the types of the arguments that follow it, and Danvy typed
printf in ML by changing that string's representation~\cite{danvy1998unparsing}. Before strings could be constant expressions, Porkol\'ab and Sinkovics
built a parser generator from template metaprograms, with a type-safe printf as its
demonstration~\cite{porkolab2010dsl}. Class types as template parameters let a string literal be a
template argument~\cite{snyder2018p0732}. Compile-time regular expressions pass a pattern as such an
argument and check it when the program compiles~\cite{dusikova2019p1433,dusikova2026ctre}. The
standard formatting library makes an invalid format string
ill-formed~\cite{zverovich2021p2216,zverovich2026fmt}. A Scala library checks regular expressions
during type checking, with match types~\cite{blanvillain2022regex}. Constant evaluation itself has grown from generalised constant
expressions~\cite{dosreis2010constexpr,stroustrup2020hopl}. The standard leaves the limits of constant evaluation to the
implementation~\cite{smith2020n4861}. Clang and MSVC expose them as a number of steps, set with
\ClangSteps\ and \MsvcSteps~\cite{llvm2026clangmanual,microsoft2026constexpr}. A table built at compile
time should cost nothing when the program runs, which is C++'s aim of zero-overhead
abstraction~\cite{stroustrup2012foundations}. H7 tests that aim on v2's tables.

\subsection{Measurement}

Small changes to an experimental setup can bias a performance evaluation~\cite{mytkowicz2009wrong}.
Benchmarks need repetition at several levels and honest estimates of variation
\cite{georges2007rigorous,kalibera2013rigorous,hoefler2015scientific,curtsinger2013stabilizer,laaber2019cloud}.
Counts that should be exact, such as retired instructions, vary from run to
run~\cite{weaver2013nondeterminism}, so we report counters as medians over processes. The harness
reads counters through the Linux interface~\cite{manpages2026perfeventopen,browne2000papi} and times
with nanobench~\cite{leitnerankerl2026nanobench}. The intervals are bias-corrected and accelerated
bootstrap intervals~\cite{efron1987bca,efron1993bootstrap,diciccio1996bootstrap}, cross-checked
against SciPy's~\cite{virtanen2020scipy}. The family of cells is controlled with Holm's step-down
procedure~\cite{holm1979simple}. The tests against a margin follow the framework of non-inferiority
and equivalence testing~\cite{schuirmann1987tost,wellek2010testing}.

\subsection{What This Paper Adds}

Typed frameworks check routes when the program compiles, and C++ libraries check strings the same
way. The sources we read report no comparison of lookup time against radix-tree routers, for typed
frameworks or for C++ string checks. Fast routers document radix trees that match without garbage or
copies~\cite{schmidt2019httprouter,ahmed2026matchit}. None of the C++ sources we read checks a route's
pattern, its handler's signature and its duplicates together when the program compiles. Yesod does
all three in Haskell, and reports no comparison of lookup time. The router evaluations we found are
benchmark suites kept beside routers, of Go routers, of Rust routers and of API gateways on real route
sets~\cite{schmidt2020gohttproutingbench,ahmed2026matchit,sworn2026wayfind,vm0012024gateways}. One
more compares a router with one table-based router~\cite{zhang2015urlrouting}. We found none that
measures routers of several languages on the same tables, checks their answers, and decides
under pre-specified tests. This paper does all three for one C++ route matcher. We do not claim that
v2 is the first typed router, or the first segment trie with hashed children. Pistache's tree and the
router of Zhang and Xu are both segment structures with hashed literal children.
